\chapter{緒論}\label{chap:introduction}
\section{研究背景}
【佔位資料】本論文示例僅用來檢查排版，不是可供引用的研究成果。
中文內文採全形括號（例如：示意資料），英文文字使用 Times New Roman。
正文段落採左右對齊；每段首行縮排兩個中文字，章節層級依序顯示。
本段刻意保留足夠長度，使自動換行、行距與左右邊界可在輸出檔中量測。

第二段示範段落間不額外加空白行。研究問題、研究對象與研究方法均須由作者
填入實際內容；本模板不預設特定研究設計，也不要求採用範例的章節名稱。
圖~\ref{fig:workflow} 及表~\ref{tab:sample} 的標題、註解與來源可供格式檢查。

\section{研究目的}
本示例驗證資料輸入、分析及呈現的關係。請以自己的研究目的替換本段。
本文的其他版型見第~\ref{chap:method} 章；附錄包含兩組獨立公式編號。

\subsection{研究範圍}
此小節示範第三層標題及內文。所有數值都是虛構測試資料；
表內括號採半形，如 sample（n），不應誤認為真實觀察結果。

\begin{figure}[htbp]
  \centering
  \setlength{\fboxsep}{10bp}
  \fbox{\begin{minipage}[c][25mm][c]{\dimexpr\linewidth-2\fboxsep-2\fboxrule\relax}
    \centering 佔位流程圖：資料輸入 \(\longrightarrow\) 分析 \(\longrightarrow\) 呈現
  \end{minipage}}
  \caption{研究流程示意圖（佔位資料）}\label{fig:workflow}
  \ccunote{方框為可替換的圖形區域，正式論文可改用外部圖檔。}
  \ccusource{本研究自行繪製（佔位資料）。}
\end{figure}
